%==================================================================
% Ini adalah bab 2
% Silahkan edit sesuai kebutuhan, baik menambah atau mengurangi \section, \subsection
%==================================================================

\chapter[KAJIAN PUSTAKA]{\\ KAJIAN PUSTAKA}

\section{Kajian Teori}

\subsection{INFINITE}

Divisi Teknologi Informasi INFINITE atau lebih sering disebut sebagai INFINITE, adalah satu dari empat divisi di bawah naungan Unit Kegiatan Mahasiswa (UKM) Rekayasa Teknologi (Restek) Universitas Negeri Yogyakarta (UNY). Organisasi ini telah ada sejak tahun 2014, dan berfungsi sebagai wadah bagi mahasiswa untuk mengembangkan diri di bidang teknologi informasi, baik melalui kegiatan pengembangan keterampilan teknis, kompetisi, maupun proyek kolaboratif. INFINITE berperan penting dalam memfasilitasi pengembangan pengetahuan dan kemampuan anggota dan membangun jejaring profesional di kalangan mahasiswa yang memiliki minat dalam teknologi informasi \citep{INFINITE_2026_Website}.

Dalam operasional hariannya, INFINITE mengelola beragam informasi yang kompleks melalui berbagai sistem aplikasi yang dimiliki. Salah satu dari sistem tersebut adalah sistem informasi manajemen organisasi utama yang menjadi pusat pengelolaan data inti organisasi. Sistem ini mencakup berbagai data seperti: (1)~data anggota, meliputi: profil akademik personal yaitu nama, Nomor Induk Mahasiswa (NIM), strata, fakultas, dan program studi; kontak yaitu nomor telepon dan email; serta status keanggotaan yaitu \textit{role} atau persona, tanggal mulai dan berakhir, status aktif keanggotaan, dan status istimewa keanggotaan; (2)~data pengurus berupa \textit{Core Team} dan \textit{Community Group Administrator}, meliputi: struktur pengurus berupa divisi; dan profil pengurus berupa jabatan, periode masa jabatan, serta foto pengurus; (3)~data kompetisi, meliputi: informasi kompetisi yaitu nama kompetisi, penyelenggara, tipe penyelenggara, tingkat kompetisi, dan cabang kompetisi; tim-tim yang dibentuk oleh anggota untuk mengikuti kompetisi yaitu nama tim, nama ketua, dan nama-nama anggota; proposal pengajuan pendanaan oleh tim yaitu detail kompetisi, tanggal pelaksanaan kompetisi, dokumen bukti lolos kompetisi, dan dokumen proposal pendanaan; serta dokumentasi prestasi yang diraih yaitu detail kompetisi, tanggal pelaksanaan kompetisi, dokumen sertifikat, dan dokumen foto kegiatan; serta (4)~data testimoni, baik dari alumni tentang dampak INFINITE terhadap karier yaitu nama alumni, posisi pekerjaan, serta testimoni; maupun testimoni proyek yang pernah dikerjakan oleh anggota yaitu nama proyek, deskripsi proyek, dan dokumentasi foto proyek.

\subsection{Sistem Informasi}

Sistem informasi merupakan kombinasi terorganisir dari manusia, perangkat keras, perangkat lunak, jaringan komunikasi, sumber data, serta kebijakan dan prosedur yang menyimpan, mengambil, mengubah, dan menyebarkan informasi dalam suatu organisasi \citep{OBrien_2012_MIS}. Menurut \citet{Laudon_2020_MIS}, sistem informasi dapat didefinisikan secara teknis sebagai sekumpulan komponen yang saling terkait yang mengumpulkan, memproses, menyimpan, dan mendistribusikan informasi untuk mendukung pengambilan keputusan, koordinasi, dan pengendalian dalam organisasi. Sistem informasi juga membantu manajer dan pekerja dalam menganalisis masalah, memvisualisasikan subjek yang kompleks, dan menciptakan produk baru.

Komponen utama sistem informasi terdiri dari lima elemen fundamental yang saling berinteraksi. Pertama, \textit{hardware} mencakup perangkat fisik seperti komputer, server, dan perangkat jaringan yang digunakan untuk memproses dan menyimpan data. Kedua, \textit{software} meliputi program aplikasi dan sistem operasi yang mengendalikan operasi perangkat keras dan memungkinkan pemrosesan data. Ketiga, data merupakan fakta mentah yang diorganisasi dan diproses menjadi informasi yang bermakna bagi pengguna. Keempat, prosedur adalah serangkaian instruksi dan aturan yang mengatur bagaimana komponen lain bekerja sama untuk memproses data menjadi informasi. Kelima, manusia (\textit{people}) mencakup pengguna akhir, administrator sistem, dan profesional teknologi informasi yang mengoperasikan dan mengelola sistem \citep{OBrien_2012_MIS}.

\subsection{\textit{Backend} dengan Laravel}

Laravel adalah kerangka kerja (\textit{framework}) aplikasi web berbasis PHP yang dikembangkan oleh Taylor Otwell pada tahun 2011. Laravel dirancang dengan filosofi pengembangan yang elegan dan ekspresif, menyediakan sintaks yang bersih serta fitur-fitur modern untuk mempercepat pengembangan aplikasi web \citep{Laravel_2026_Docs}. Sebagai salah satu kerangka kerja PHP yang paling populer, Laravel menawarkan berbagai fitur bawaan seperti sistem \textit{routing} yang fleksibel, \textit{Object-Relational Mapping} (ORM) melalui Eloquent, sistem migrasi basis data, antrian pekerjaan (\textit{job queues}), serta sistem autentikasi dan otorisasi yang komprehensif. Studi komparatif terhadap kerangka kerja PHP menunjukkan bahwa Laravel memiliki keunggulan dalam hal dokumentasi, komunitas pengembang, dan ekosistem paket pendukung \citep{Laaziri_2019_PHPFrameworks}.

\subsubsection{PHP}

PHP adalah bahasa pemrograman skrip sisi server yang awalnya dikembangkan oleh Rasmus Lerdorf pada tahun 1994 dengan nama \textit{Personal Home Page Tools}. Seiring perkembangannya, PHP bertransformasi menjadi bahasa pemrograman yang lengkap dan namanya diubah menjadi akronim rekursif \textit{PHP: Hypertext Preprocessor} \citep{PHP_2026_History}. PHP dirancang khusus untuk pengembangan web dan dapat disisipkan ke dalam HTML, menjadikannya pilihan populer untuk membangun aplikasi web dinamis.

Dibandingkan dengan bahasa pemrograman \textit{backend} lainnya seperti Python, Java, atau Node.js, PHP memiliki beberapa keunggulan. Pertama, PHP memiliki kurva pembelajaran yang relatif landai dan sintaks yang mudah dipahami. Kedua, ekosistem PHP sangat matang dengan ketersediaan banyak kerangka kerja seperti Laravel, Symfony, dan CodeIgniter. Ketiga, sebagian besar penyedia \textit{hosting} web mendukung PHP secara \textit{native}, memudahkan proses \textit{deployment}. Namun demikian, PHP juga memiliki keterbatasan seperti inkonsistensi penamaan fungsi bawaan dan penanganan \textit{concurrency} yang kurang optimal dibandingkan Node.js. Dalam konteks proyek ini, PHP dipilih untuk menjaga kemudahan transisi dari sistem informasi versi sebelumnya, dan karena kerangka kerja Laravel yang digunakan sudah menyediakan abstraksi yang elegan untuk pengembangan API \textit{backend} serta memiliki dokumentasi dan komunitas yang sangat aktif.

\subsubsection{Arsitektur \textit{Model-View-Controller} (MVC)}

MVC adalah pola arsitektur perangkat lunak yang memisahkan aplikasi menjadi tiga komponen utama yang saling terhubung namun memiliki tanggung jawab yang berbeda. Arsitektur ini pertama kali diperkenalkan oleh Trygve Reenskaug pada tahun 1979 di Xerox PARC dalam konteks pengembangan antarmuka pengguna Smalltalk-80. Konsep MVC kemudian dipopulerkan dan didokumentasikan secara lebih luas oleh \citet{Krasner_1988_MVC} dalam publikasi yang menjadi referensi klasik untuk implementasi arsitektur ini. Pemisahan ini bertujuan untuk meningkatkan modularitas, memudahkan pemeliharaan, dan memungkinkan pengembangan paralel oleh tim yang berbeda. Dalam konteks teknologi web modern, arsitektur MVC telah berevolusi dengan berbagai adaptasi untuk memenuhi kebutuhan aplikasi web yang kompleks \citep{Sastypratiwi_2019_MVCComponentWeb}.

Komponen pertama, \textit{Model}, bertanggung jawab atas logika bisnis dan pengelolaan data aplikasi. \textit{Model} merepresentasikan struktur data, aturan validasi, dan operasi terhadap basis data. Dalam Laravel, \textit{Model} diimplementasikan melalui Eloquent ORM yang menyediakan antarmuka objek untuk berinteraksi dengan tabel basis data. Komponen kedua, \textit{View}, menangani presentasi data kepada pengguna. \textit{View} menerima data dari \textit{Controller} dan merendernya dalam format yang sesuai untuk ditampilkan. Komponen ketiga, \textit{Controller}, berperan sebagai perantara antara \textit{Model} dan \textit{View}. \textit{Controller} menerima permintaan dari pengguna, memproses logika yang diperlukan dengan memanggil \textit{Model}, dan mengembalikan respons melalui \textit{View} \citep{Laravel_2026_Docs}.

Dalam implementasi Laravel \textit{headless}, arsitektur MVC diadaptasi tanpa komponen \textit{Blade templating} untuk menampilkan HTML. Sebagai gantinya, Laravel menggunakan \textit{Resources} dan \textit{Collections} untuk memformat data \textit{Model} menjadi respons JSON yang konsisten. Laravel menyediakan fitur-fitur seperti \textit{route model binding}, \textit{middleware} untuk autentikasi dan otorisasi, serta \textit{resource controllers} yang membantu untuk pengembangan API \citep{Laravel_2026_Docs}.


\subsubsection{\textit{RESTful Application Programming Interface} (API)}

\textit{Representational State Transfer} (REST) adalah gaya arsitektur untuk sistem terdistribusi yang pertama kali diperkenalkan oleh Roy Fielding dalam disertasinya pada tahun 2000 \citep{Fielding_2000_REST}. REST mendefinisikan sekumpulan batasan arsitektural yang, ketika diterapkan secara keseluruhan, menekankan skalabilitas interaksi antar komponen, antarmuka yang seragam, dan \textit{deployment} komponen yang independen. API yang mengikuti prinsip REST disebut sebagai \textit{RESTful} API.

Prinsip utama REST meliputi: (1)~arsitektur \textit{client-server} yang memisahkan antarmuka pengguna dari penyimpanan data; (2)~\textit{statelessness}, di mana setiap permintaan dari klien harus berisi semua informasi yang diperlukan untuk memproses permintaan; (3)~\textit{cacheability} yang memungkinkan respons ditandai sebagai dapat di-\textit{cache} atau tidak; (4)~sistem berlapis (\textit{layered system}) yang memungkinkan penambahan lapisan perantara; dan (5)~antarmuka yang seragam menggunakan metode HTTP standar (GET, POST, PUT, DELETE) untuk operasi \textit{Create, Read, Update, Delete} (CRUD) \citep{Richardson_2013_RESTfulWebAPIs}.

\subsubsection{Basis Data PostgreSQL}

PostgreSQL adalah sistem manajemen basis data relasional objek (\textit{Object-Relational Database Management System}/ORDBMS) sumber terbuka yang dikembangkan sejak tahun 1986 sebagai bagian dari proyek POSTGRES di University of California, Berkeley \citep{PostgreSQL_2026_About}. PostgreSQL dikenal karena keandalan, integritas data, dan kepatuhannya terhadap standar SQL. Sebagai basis data sumber terbuka, PostgreSQL dapat digunakan, dimodifikasi, dan didistribusikan secara bebas tanpa biaya lisensi.

Fitur-fitur utama PostgreSQL meliputi: (1)~kepatuhan penuh terhadap properti \textit{Atomicity, Consistency, Isolation, Durability} (ACID) yang menjamin integritas transaksi; (2)~dukungan tipe data yang ekstensif termasuk JSON/JSONB untuk data semi-terstruktur; (3)~sistem indeks yang canggih termasuk B-tree, Hash, GiST, dan GIN; (4)~\textit{full-text search} bawaan untuk pencarian teks; serta (5)~kemampuan ekstensi yang memungkinkan penambahan fungsi kustom \citep{Ferrari_2023_PostgreSQL}. Dalam konteks proyek INFINITE, PostgreSQL dipilih karena beberapa alasan: pertama, PostgreSQL telah digunakan sebelumnya untuk beberapa aplikasi internal INFINITE sehingga tidak memerlukan \textit{deployment} baru; kedua, sebagai basis data sumber terbuka, PostgreSQL tidak memerlukan biaya lisensi yang sejalan dengan prinsip pengembangan organisasi mahasiswa; dan ketiga, PostgreSQL memiliki performa yang sangat baik untuk beban kerja transaksional maupun analitis.

\subsection{\textit{Frontend} dengan Next.js}

Next.js adalah kerangka kerja React yang dikembangkan oleh Vercel untuk membangun aplikasi web dengan kemampuan \textit{rendering} yang fleksibel \citep{Nextjs_2026_Docs}. Next.js memperluas kemampuan React dengan menyediakan fitur-fitur seperti \textit{Server-Side Rendering} (SSR), \textit{Static Site Generation} (SSG), \textit{Incremental Static Regeneration} (ISR), dan sistem \textit{routing} berbasis \textit{file system}. Kombinasi fitur-fitur ini memungkinkan pengembang untuk mengoptimalkan performa aplikasi sesuai dengan kebutuhan masing-masing halaman.

React.js sendiri adalah pustaka JavaScript untuk membangun antarmuka pengguna yang dikembangkan oleh Meta (sebelumnya Facebook) dan dirilis sebagai sumber terbuka pada tahun 2013 \citep{React_2026_Docs}. React menggunakan pendekatan berbasis komponen (\textit{component-based}) yang memungkinkan pengembang membangun UI dari potongan-potongan kode yang terisolasi dan dapat digunakan kembali \citep{Kaushalya_2021_ReactComponentMigration}. Konsep utama React meliputi: (1)~\textit{Virtual DOM} yang mengoptimalkan pembaruan UI dengan meminimalkan manipulasi DOM langsung; (2)~\textit{JavaScript XML} (JSX) yang memungkinkan penulisan markup dalam kode JavaScript; (3)~\textit{unidirectional data flow} yang membuat aliran data lebih mudah diprediksi; serta (4)~\textit{hooks} yang memungkinkan penggunaan \textit{state} dan fitur React lainnya dalam komponen fungsional.

\subsubsection{TypeScript}

TypeScript adalah bahasa pemrograman sumber terbuka yang dikembangkan dan dipelihara oleh Microsoft sejak tahun 2012 \citep{TypeScript_2026_Docs}. TypeScript merupakan \textit{superset} dari JavaScript yang menambahkan sistem tipe statis opsional dan fitur-fitur pemrograman berorientasi objek yang lebih kaya. Kode TypeScript dikompilasi (\textit{transpiled}) menjadi JavaScript murni yang dapat dijalankan di berbagai lingkungan seperti peramban, Node.js, atau \textit{runtime} JavaScript lainnya.

Sistem tipe TypeScript memungkinkan pendeteksian kesalahan pada saat kompilasi alih-alih saat eksekusi, yang secara signifikan meningkatkan keandalan kode \citep{Bierman_2014_TypeScript}. Fitur-fitur utama TypeScript meliputi: (1)~anotasi tipe untuk variabel, parameter fungsi, dan nilai kembalian; (2)~\textit{interfaces} dan \textit{type aliases} untuk mendefinisikan bentuk objek; (3)~\textit{generics} untuk membuat komponen yang dapat bekerja dengan berbagai tipe; (4)~\textit{enums} untuk mendefinisikan sekumpulan konstanta bernama; serta (5)~dukungan \textit{IntelliSense} yang lebih baik di \textit{Integrated Development Environment} (IDE). Dalam pengembangan aplikasi \textit{frontend} dengan Next.js, penggunaan TypeScript meningkatkan \textit{maintainability} kode, memudahkan \textit{refactoring}, dan memperbaiki kolaborasi tim melalui dokumentasi tipe yang eksplisit.

\subsubsection{\textit{Clean Architecture}}

\textit{Clean Architecture} adalah pendekatan arsitektur perangkat lunak yang diperkenalkan oleh Robert C. Martin (Uncle Bob) yang menekankan pemisahan kepentingan (\textit{separation of concerns}) dan independensi dari kerangka kerja, basis data, serta antarmuka pengguna \citep{Martin_2017_CleanArchitecture}. Prinsip utama \textit{Clean Architecture} adalah \textit{Dependency Rule} yang menyatakan bahwa dependensi kode sumber hanya boleh mengarah ke dalam, menuju kebijakan tingkat tinggi (\textit{high-level policies}).

Arsitektur ini mengorganisasi kode ke dalam beberapa lapisan konsentris dengan tanggung jawab yang berbeda. Lapisan terdalam adalah \textit{Domain} (atau \textit{Entities}) yang berisi logika bisnis inti dan aturan bisnis yang paling umum. Lapisan \textit{Application} (atau \textit{Use Cases}) berisi logika aplikasi spesifik yang mengorkestrasi aliran data ke dan dari entitas. Lapisan \textit{Infrastructure} (atau \textit{Interface Adapters}) berisi implementasi konkret dari antarmuka yang didefinisikan di lapisan dalam, seperti repositori basis data atau klien API. Lapisan terluar adalah \textit{Presentation} (atau \textit{Frameworks and Drivers}) yang berisi detail implementasi seperti kerangka kerja UI atau pustaka eksternal.

Dalam penelitian ini, \textit{Clean Architecture} diterapkan melalui struktur direktori yang mencerminkan lapisan-lapisan tersebut: direktori \texttt{domain} untuk entitas dan aturan bisnis, \texttt{application} untuk \textit{use cases} dan layanan aplikasi, \texttt{infrastructure} untuk implementasi repositori dan integrasi API, serta \texttt{presentation} untuk komponen React dan halaman Next.js. Pemisahan ini meningkatkan modularitas dan memungkinkan perubahan pada satu lapisan tanpa memengaruhi lapisan lainnya.

\subsection{\textit{Deployment} dengan Kubernetes dan GitOps}

\subsubsection{Komputasi Awan}

Dalam operasional hariannya, INFINITE telah memanfaatkan komputasi awan untuk menjalankan berbagai sistem aplikasi yang dimiliki. Menurut National Institute of Standards and Technology (NIST), komputasi awan atau \textit{cloud computing} adalah suatu model yang memungkinkan akses secara \textit{on-demand} melalui jaringan kepada sekumpulan sumber daya komputasi bersama yang dapat disediakan dan dilepaskan dengan cepat. Karakteristik utama dari komputasi awan meliputi: (1)~layanan mandiri sesuai kebutuhan, pengguna dapat mengalokasikan sumber daya komputasi secara terprogram tanpa interaksi manusia; (2)~akses jaringan luas, tersedia melalui jaringan standar yang mendukung berbagai platform pengguna; (3)~penggabungan sumber daya, pengelolaan sumber daya dilakukan secara terpusat dengan model \textit{multi-tenant}; (4)~elastisitas cepat, memungkinkan penyediaan dan pelepasan sumber daya secara dinamis sesuai beban kerja; dan (5)~layanan terukur, optimasi penggunaan sumber daya melalui mekanisme pengukuran yang transparan \citep{Mell_2011_CloudComputing}. Secara umum, layanan komputasi awan terbagi ke dalam 3 model, yaitu \textit{Infrastructure as a Service} (IaaS), \textit{Platform as a Service} (PaaS), dan \textit{Software as a Service} (SaaS). IaaS menyediakan sumber daya komputasi dasar seperti mesin virtual, penyimpanan, dan jaringan; PaaS menyediakan platform untuk pengembangan aplikasi yang mencakup sistem operasi, \textit{runtime}, dan alat pengembangan; sedangkan SaaS menyediakan aplikasi siap pakai yang dapat diakses melalui web tanpa perlu instalasi lokal \citep{Younis_2024_Comprehensive}.

\subsubsection{Kontainer dan Kubernetes}

Seiring dengan adopsi komputasi awan yang meluas, tantangan utama yang dihadapi organisasi bergeser dari penyediaan infrastruktur ke pengelolaan kompleksitas sistem terdistribusi. Untuk mengatasi tantangan ini, platform orkestrasi seperti Kubernetes telah menjadi fondasi penting dalam membangun dan mengelola aplikasi \textit{cloud-native} yang skalabel dan dapat diandalkan. Kubernetes adalah platform orkestrasi kontainer sumber terbuka yang awalnya dikembangkan oleh Google berdasarkan pengalaman mengelola sistem berskala besar melalui proyek internal Borg \citep{Burns_2016_Borg}. Platform ini menyediakan abstraksi deklaratif untuk mengelola siklus hidup aplikasi terkontainerisasi, termasuk penjadwalan otomatis, penskalaan horizontal, pemulihan mandiri, penemuan layanan, serta penyeimbangan beban. Kubernetes mengorganisasi sumber daya dalam unit-unit seperti \textit{Pod} (unit terkecil yang dapat dijadwalkan), \textit{Deployment} (manajemen replika dan pembaruan), \textit{Service} (abstraksi jaringan), dan \textit{Namespace} (isolasi logis) \citep{Burns_2019_Kubernetes}. Hingga kini, adopsi Kubernetes telah meningkat sangat pesat. Survei Cloud Native Computing Foundation (CNCF) pada tahun 2025 menunjukkan bahwa 82\% organisasi telah menggunakan Kubernetes dalam lingkungan produksi \citep{CNCF_2026_Kubernetes}.

\subsubsection{\textit{Continuous Integration/Continuous Deployment} (CI/CD)}

\textit{Continuous Integration} (CI) dan \textit{Continuous Deployment} (CD) adalah praktik pengembangan perangkat lunak yang bertujuan untuk mengotomatisasi dan mempercepat siklus rilis perangkat lunak \citep{Shahin_2017_CICD}. CI mengacu pada praktik mengintegrasikan perubahan kode dari beberapa kontributor ke dalam repositori bersama secara sering, biasanya beberapa kali sehari, disertai dengan \textit{build} dan pengujian otomatis untuk mendeteksi kesalahan sedini mungkin. CD memperluas CI dengan mengotomatisasi proses \textit{deployment} ke lingkungan produksi atau \textit{staging}. Terdapat dua varian CD: \textit{Continuous Delivery} yang mengotomatisasi proses hingga siap di-\textit{deploy} namun memerlukan persetujuan manual, dan \textit{Continuous Deployment} yang sepenuhnya otomatis hingga produksi \citep{HumbleFarley_2010_CD}.

\subsubsection{GitHub Actions}

GitHub Actions adalah platform CI/CD yang terintegrasi langsung dengan GitHub, memungkinkan otomatisasi alur kerja pengembangan perangkat lunak \citep{GitHubActions_2024_Docs}. Alur kerja (\textit{workflow}) didefinisikan dalam berkas YAML yang disimpan di direktori \texttt{.github/workflows/} dan dapat dipicu oleh berbagai \textit{event} seperti \textit{push}, \textit{pull request}, atau jadwal \textit{cron}. Dalam proyek ini, GitHub Actions digunakan untuk membangun \textit{container image} dari \textit{backend} Laravel dan \textit{frontend} Next.js setiap kali ada perubahan kode yang di-\textit{push} ke cabang repositori produksi.

\subsubsection{GitOps dengan ArgoCD}

GitOps adalah paradigma operasional yang menggunakan Git sebagai sumber kebenaran tunggal (\textit{single source of truth}) untuk infrastruktur deklaratif dan aplikasi. Konsep GitOps merepresentasikan evolusi praktik DevOps yang memanfaatkan sistem kontrol versi sebagai mekanisme sentral untuk manajemen infrastruktur dan \textit{deployment} aplikasi \citep{Beetz_2022_GitOps}. Prinsip utama GitOps meliputi: (1)~konfigurasi sistem yang dideskripsikan secara deklaratif; (2)~status sistem yang diinginkan disimpan di Git; (3)~perubahan yang disetujui diterapkan secara otomatis ke sistem; dan (4)~agen perangkat lunak memastikan kebenaran dan melaporkan penyimpangan. Pendekatan GitOps telah menjadi bagian penting dari ekosistem DevOps modern, memberikan visibilitas, auditabilitas, dan konsistensi dalam pengelolaan infrastruktur \citep{Leite_2019_GitOps}. ArgoCD adalah salah satu contoh alat GitOps sumber terbuka untuk Kubernetes yang secara kontinu memantau repositori Git dan menyinkronkan status klaster Kubernetes agar sesuai dengan konfigurasi yang didefinisikan \citep{ArgoCD_2024_Docs}. ArgoCD menyediakan fitur seperti deteksi \textit{drift}, \textit{rollback} otomatis, serta antarmuka web untuk visualisasi dan manajemen aplikasi.

\subsection{Autentikasi dan \textit{Single Sign-On} (SSO)}

Autentikasi adalah proses verifikasi identitas pengguna atau sistem yang mengakses sumber daya digital \citep{Grassi_2017_NIST800}. Dalam konteks aplikasi web, autentikasi memastikan bahwa pengguna yang mengakses sistem adalah benar-benar pihak yang mengklaim identitasnya. Menurut \citet{Ometov_2018_MultiFactorAuth}, autentikasi dapat diklasifikasikan berdasarkan faktor yang digunakan: (1)~\textit{knowledge-based} seperti kata sandi atau PIN, (2)~\textit{possession-based} seperti kartu atau perangkat \textit{hardware token}, dan (3)~\textit{inherence-based} seperti sidik jari atau pengenalan wajah. Dalam aplikasi modern, kombinasi beberapa faktor (\textit{multi-factor authentication}) semakin umum digunakan untuk meningkatkan keamanan.

SSO adalah mekanisme autentikasi yang memungkinkan pengguna untuk mengakses berbagai aplikasi dengan menggunakan satu set kredensial tunggal \citep{Pashalidis_2003_SSO}. Implementasi SSO memberikan beberapa manfaat signifikan. Pertama, SSO meningkatkan pengalaman pengguna dengan mengeliminasi kebutuhan untuk mengingat dan mengelola kredensial yang berbeda untuk setiap aplikasi. Kedua, SSO mengurangi beban administratif dalam pengelolaan akun pengguna dan pengaturan ulang kata sandi. Ketiga, SSO dapat meningkatkan keamanan dengan memungkinkan penerapan kebijakan autentikasi yang konsisten di seluruh aplikasi \citep{Liu_2018_AndroidSSOSecurity}. Dalam konteks organisasi yang memiliki ekosistem aplikasi yang kompleks seperti INFINITE, SSO menjadi solusi yang efektif untuk menyederhanakan manajemen identitas dan akses.

\subsubsection{OAuth 2.0}

OAuth 2.0 adalah kerangka kerja autorisasi yang memungkinkan aplikasi pihak ketiga untuk mendapatkan akses terbatas ke suatu layanan atas nama pemilik sumber daya \citep{Hardt_2012_OAuth2}. Berbeda dengan protokol autentikasi, OAuth 2.0 dirancang khusus untuk skenario autorisasi di mana aplikasi memerlukan akses ke sumber daya tanpa perlu mengetahui kredensial pengguna. Menurut \citet{Fett_2016_OAuth2Security}, OAuth 2.0 mendefinisikan empat peran utama: (1)~\textit{Resource Owner}, pemilik sumber daya, biasanya pengguna; (2)~\textit{Resource Server}, server yang menyimpan sumber daya yang dilindungi; (3)~\textit{Client}, aplikasi yang meminta akses; dan (4)~\textit{Authorization Server}, server yang mengeluarkan \textit{access token} setelah berhasil mengautentikasi \textit{resource owner}.

OAuth 2.0 mendefinisikan beberapa \textit{grant type} untuk berbagai kasus penggunaan. \textit{Authorization Code Grant} adalah yang paling aman dan umum digunakan untuk aplikasi web, di mana \textit{client} menerima kode otorisasi yang kemudian ditukar dengan \textit{access token}. \textit{Implicit Grant} dirancang untuk aplikasi berbasis peramban namun tidak lagi direkomendasikan karena masalah keamanan. \textit{Resource Owner Password Credentials Grant} memungkinkan \textit{client} menggunakan kredensial pengguna secara langsung, namun hanya cocok untuk aplikasi yang sangat tepercaya. \textit{Client Credentials Grant} digunakan untuk autentikasi \textit{machine-to-machine} tanpa konteks pengguna manusia \citep{Lodderstedt_2020_OAuth2BestPractices}.

\subsubsection{OpenID Connect (OIDC)}

OIDC adalah lapisan identitas yang dibangun di atas OAuth 2.0 yang menambahkan kemampuan autentikasi standar \citep{Sakimura_2014_OpenIDConnect}. Sementara OAuth 2.0 berfokus pada autorisasi, OIDC menyediakan mekanisme untuk memverifikasi identitas pengguna dan mendapatkan informasi profil dasar. OIDC menambahkan konsep \textit{ID Token}, yang merupakan \textit{JSON Web Token} (JWT) yang berisi klaim tentang autentikasi pengguna dan atribut identitas \citep{Jones_2015_JWT}. Komponen utama OIDC mencakup: (1)~\textit{OpenID Provider} (OP), yaitu server OAuth 2.0 yang mampu mengautentikasi pengguna dan menyediakan klaim tentang autentikasi tersebut; (2)~\textit{Relying Party} (RP), yaitu aplikasi yang menggunakan OIDC untuk autentikasi; dan (3)~\textit{End User}, yaitu pengguna yang identitasnya diverifikasi.

OIDC mendefinisikan beberapa \textit{endpoint} standar yang harus disediakan oleh \textit{OpenID Provider}: (1)~\textit{Authorization endpoint} untuk memulai proses autentikasi, (2)~\textit{Token endpoint} untuk menukar kode otorisasi dengan token, (3)~\textit{UserInfo endpoint} untuk mendapatkan klaim tambahan tentang pengguna, dan (4)~\textit{Discovery endpoint} yang menyediakan metadata konfigurasi OIDC secara otomatis. Standardisasi ini memungkinkan interoperabilitas yang tinggi antara berbagai implementasi. Namun, implementasi OIDC harus memperhatikan aspek keamanan secara cermat, termasuk validasi token yang benar, penanganan alur autentikasi yang aman, dan perlindungan terhadap serangan seperti \textit{token substitution} dan \textit{session fixation} \citep{Mladenov_2016_OIDCSecurity,Fett_2014_OIDCFormal}.

OIDC juga mendefinisikan tiga alur autentikasi utama dengan karakteristik yang berbeda. Pertama, \textit{Authorization Code Flow} merupakan alur yang paling aman di mana klien menerima kode otorisasi dari \textit{authorization endpoint}, kemudian menukarkan kode tersebut dengan \textit{ID token} dan \textit{access token} melalui \textit{Token endpoint}. Alur ini cocok untuk aplikasi web tradisional yang memiliki komponen \textit{server-side} karena \textit{client secret} dapat disimpan dengan aman di server. Kedua, \textit{Implicit Flow} dirancang untuk aplikasi yang berjalan sepenuhnya di peramban tanpa komponen \textit{backend}, di mana token langsung dikembalikan dari \textit{Authorization endpoint} tanpa tahap pertukaran kode. Namun, alur ini tidak lagi direkomendasikan karena risiko keamanan yang terkait dengan eksposur token di URL dan ketidakmampuan untuk mengautentikasi klien secara efektif. Ketiga, \textit{Hybrid Flow} menggabungkan elemen dari kedua alur sebelumnya, memungkinkan beberapa token dikembalikan dari \textit{Authorization endpoint} dan beberapa dari \textit{Token endpoint}, memberikan fleksibilitas untuk skenario yang memerlukan akses langsung ke \textit{ID token} di \textit{frontend} sambil tetap mempertahankan keamanan untuk \textit{access token}. Untuk aplikasi web modern, \textit{Authorization Code Flow} dengan \textit{Proof Key for Code Exchange} (PKCE) direkomendasikan untuk keamanan maksimal, bahkan untuk aplikasi \textit{single-page} dan \textit{mobile} yang sebelumnya menggunakan \textit{Implicit Flow} \citep{Bradley_2025_OAuth2Browsers}.

\subsubsection{\textit{Security Assertion Markup Language} (SAML)}

SAML adalah standar terbuka berbasis XML untuk pertukaran data autentikasi dan autorisasi antara \textit{identity provider} dan \textit{service provider} \citep{Hughes_2005_SAMLProfiles}. SAML telah lama digunakan dalam lingkungan \textit{enterprise} untuk mengimplementasikan SSO, terutama untuk aplikasi berbasis web tradisional. Berbeda dengan OAuth 2.0 dan OIDC yang menggunakan format JSON dan REST, SAML menggunakan format XML dan umumnya dikomunikasikan melalui HTTP POST atau redirect.

Arsitektur SAML terdiri dari tiga komponen utama: (1)~\textit{Identity Provider} (IdP) yang mengautentikasi pengguna dan mengeluarkan \textit{assertion} yang berisi informasi autentikasi dan atribut pengguna, (2)~\textit{Service Provider} (SP) yang menyediakan layanan dan menerima \textit{assertion} dari IdP untuk membuat keputusan akses, dan (3)~\textit{User Agent} (biasanya peramban) yang memfasilitasi komunikasi antara IdP dan SP. SAML mendefinisikan tiga jenis \textit{assertion}: \textit{Authentication Assertion} untuk membuktikan autentikasi pengguna, \textit{Attribute Assertion} untuk menyampaikan atribut pengguna, dan \textit{Authorization Decision Assertion} untuk keputusan akses \citep{Cantor_2008_SAMLBindings}.

Dalam alur SSO berbasis SAML, proses dimulai ketika pengguna mencoba mengakses \textit{service provider}. SP mengarahkan pengguna ke IdP untuk autentikasi. Setelah pengguna berhasil diautentikasi oleh IdP, IdP mengeluarkan \textit{SAML assertion} yang ditandatangani secara digital dan mengirimkannya kembali ke SP melalui peramban pengguna. SP memvalidasi \textit{assertion} dan memberikan akses kepada pengguna. SAML mendukung dua binding utama: HTTP Redirect Binding dan HTTP POST Binding, yang menentukan bagaimana pesan SAML ditransmisikan \citep{Mainka_2017_SoK_OpenIDConnect}.

\subsubsection{Perbandingan Protokol SSO}

Ketiga protokol SSO memiliki karakteristik dan kasus penggunaan yang berbeda. SAML lebih banyak digunakan dalam lingkungan \textit{enterprise} yang sudah mapan, terutama untuk integrasi dengan sistem \textit{legacy} dan aplikasi berbasis web tradisional. SAML menyediakan fitur yang sangat lengkap dan matang, namun kompleksitas XML dan overhead protokol dapat menjadi hambatan untuk aplikasi modern. OAuth 2.0 dan OIDC lebih populer dalam ekosistem aplikasi modern, terutama untuk aplikasi \textit{mobile}, \textit{Single-Page Application} (SPA), dan API. OAuth 2.0 cocok untuk skenario di mana aplikasi pihak ketiga memerlukan akses terbatas ke sumber daya pengguna tanpa perlu mengetahui kredensial, seperti integrasi dengan layanan media sosial atau aplikasi \textit{cloud}. OIDC melengkapi OAuth 2.0 dengan menambahkan lapisan autentikasi yang terstandarisasi, menjadikannya pilihan yang ideal untuk implementasi SSO modern yang memerlukan autentikasi pengguna dan informasi profil. Analisis praktik terbaik OAuth/OIDC untuk aplikasi modern oleh \citet{Sharif_2022_OAuthOIDCNativeBCP} menunjukkan pentingnya pemilihan alur protokol yang aman sesuai konteks aplikasi. Panduan keamanan terkini dari \citet{Lodderstedt_2025_OAuth2SecurityBCP} juga menegaskan bahwa OIDC dan OAuth 2.0 telah menjadi standar \textit{de facto} untuk autentikasi dan autorisasi berbasis web maupun \textit{mobile}.

\subsection{Rekayasa Perangkat Lunak}

\subsubsection{Perangkat Lunak Berbasis \textit{Website}}

Perangkat lunak berbasis \textit{website} atau aplikasi web adalah jenis perangkat lunak yang diakses melalui peramban dan berjalan di atas protokol HTTP/HTTPS, di mana sebagian atau seluruh logika aplikasi dieksekusi di sisi server dan antarmuka pengguna ditampilkan melalui peramban \citep{Taivalsaari_2021_WebApps}. Berbeda dengan aplikasi \textit{desktop} tradisional yang harus diinstal secara lokal di setiap perangkat, aplikasi web dapat diakses dari berbagai perangkat tanpa memerlukan instalasi khusus selain peramban. Menurut \citet{Mikkonen_2019_WebPlatformCornucopia}, evolusi teknologi web dalam dekade terakhir telah mengubah web dari platform untuk menyajikan dokumen statis menjadi platform pengembangan aplikasi yang matang dan \textit{powerful}, memungkinkan pembuatan aplikasi yang kompleks dengan pengalaman pengguna yang setara atau bahkan melampaui aplikasi \textit{native}.

Keunggulan perangkat lunak berbasis \textit{website} mencakup beberapa aspek signifikan. Pertama, aksesibilitas lintas platform yang memungkinkan pengguna mengakses aplikasi dari berbagai sistem operasi (Windows, macOS, Linux) dan perangkat (\textit{desktop}, \textit{tablet}, \textit{smartphone}) selama memiliki peramban dan koneksi internet. Kedua, kemudahan distribusi dan pembaruan, di mana pengembang dapat melakukan \textit{deployment} pembaruan secara terpusat di server tanpa memerlukan pengguna untuk mengunduh dan menginstal versi baru. Ketiga, biaya pengembangan yang relatif lebih rendah karena hanya perlu mengembangkan satu basis kode untuk berbagai platform. Keempat, manajemen data terpusat yang memudahkan sinkronisasi dan \textit{backup} data pengguna \citep{Charland_2011_MobileWeb}.

Arsitektur aplikasi web modern dapat diklasifikasikan ke dalam beberapa model berdasarkan pembagian tanggung jawab antara \textit{client} dan \textit{server}. Arsitektur \textit{traditional server-side rendering} merender seluruh halaman HTML di sisi server dan mengirimkannya ke \textit{client}, memberikan keunggulan dalam hal \textit{Search Engine Optimization} (SEO) dan waktu muat awal yang cepat, namun memiliki keterbatasan dalam hal interaktivitas. Arsitektur SPA memuat seluruh aplikasi sekali di awal dan kemudian berkomunikasi dengan server melalui API untuk data, memberikan pengalaman pengguna yang lebih responsif namun dapat menghadapi tantangan dalam SEO dan waktu muat awal. Arsitektur \textit{Progressive Web App} (PWA) menggabungkan fitur aplikasi web dengan kemampuan aplikasi \textit{native} seperti kerja \textit{offline} dan notifikasi \textit{push} \citep{Biorn_2018_PWA}. Dalam perkembangan terkini, arsitektur \textit{JavaScript, APIs, and Markup} (JAMstack) telah muncul sebagai pendekatan modern yang memisahkan \textit{frontend} dari \textit{backend}, memanfaatkan API untuk komunikasi data, dan meningkatkan performa melalui \textit{pre-rendering} dan penggunaan \textit{Content Delivery Network} (CDN). Studi empiris menunjukkan bahwa JAMstack dapat meningkatkan performa aplikasi web secara signifikan dibandingkan arsitektur tradisional \citep{Markovic_2022_JAMstack}.

Dalam penelitian ini, arsitektur \textit{headless} diterapkan sebagai pendekatan untuk memisahkan \textit{backend} dan \textit{frontend} secara total. Arsitektur \textit{headless} merujuk pada sistem di mana lapisan \textit{backend} yang mengelola logika bisnis, basis data, dan API sepenuhnya terpisah dari lapisan \textit{frontend} yang mengelola presentasi dan antarmuka pengguna \citep{Hassan_2024_HeadlessCMSPerformance}. Pemisahan ini memungkinkan pengembangan, \textit{deployment}, dan penskalaan komponen secara independen. Pendekatan API-\textit{first} dalam arsitektur \textit{headless} memungkinkan desain dan implementasi API sebelum pengembangan antarmuka pengguna, memfasilitasi integrasi yang lebih baik dengan berbagai klien dan meningkatkan fleksibilitas sistem \citep{Velepucha_2023_MicroservicesSurvey}. \textit{Backend} Laravel menyediakan RESTful API yang dapat dikonsumsi oleh berbagai klien, sementara \textit{frontend} Next.js bertanggung jawab sepenuhnya untuk rendering dan interaksi pengguna. Keunggulan arsitektur \textit{headless} meliputi: (1)~fleksibilitas teknologi, di mana \textit{frontend} dan \textit{backend} dapat dikembangkan dengan teknologi yang berbeda dan diperbarui secara independen; (2)~\textit{reusability} API, di mana satu \textit{backend} dapat melayani berbagai klien seperti web, \textit{mobile}, atau IoT; (3)~peningkatan performa melalui optimasi dan \textit{caching} yang lebih baik pada setiap lapisan; dan (4)~kemudahan kolaborasi tim, di mana tim \textit{frontend} dan \textit{backend} dapat bekerja secara paralel \citep{Strapi_2026_Headless}. Namun demikian, arsitektur \textit{headless} juga memiliki tantangan seperti kompleksitas infrastruktur yang lebih tinggi, kebutuhan komunikasi API yang harus dikelola dengan baik, dan kurva pembelajaran yang lebih curam untuk pengembang yang terbiasa dengan arsitektur monolitik tradisional.

\subsubsection{Model Pengembangan Perangkat Lunak}

Model pengembangan perangkat lunak adalah kerangka kerja terstruktur yang menjelaskan proses, kegiatan, dan tahapan yang dilalui dalam membangun sistem perangkat lunak dari konsepsi awal hingga \textit{deployment} dan pemeliharaan \citep{Sommerville_2016_SE}. Model pengembangan menyediakan panduan sistematis untuk mengelola kompleksitas proyek perangkat lunak, memfasilitasi komunikasi antar anggota tim, mengalokasikan sumber daya secara efektif, dan memastikan bahwa produk akhir memenuhi kebutuhan pengguna. Menurut \citet{Pressman_2014_SE}, pemilihan model pengembangan yang tepat sangat bergantung pada karakteristik proyek seperti ukuran, kompleksitas, tingkat kepastian kebutuhan, dan sumber daya yang tersedia.

Dalam penelitian ini, model \textit{Waterfall} digunakan sebagai kerangka pengembangan sistem informasi organisasi INFINITE. Model \textit{Waterfall} pertama kali diperkenalkan oleh \citet{Royce_1970_Waterfall} dan merupakan salah satu model pengembangan perangkat lunak tertua yang mengikuti pendekatan sekuensial dan linear. Dalam model ini, setiap tahap pengembangan harus diselesaikan secara penuh sebelum melanjutkan ke tahap berikutnya, dengan aliran progres yang mengalir ke bawah layaknya air terjun. Meskipun model \textit{Waterfall} telah lama digunakan, penelitian terkini menunjukkan bahwa model ini tetap relevan untuk proyek dengan kebutuhan yang stabil dan terdefinisi dengan baik, terutama dalam konteks akademis dan organisasi dengan sumber daya terbatas \citep{Saravanos_2023_Waterfall}. Studi komparatif menunjukkan bahwa model \textit{Waterfall} cocok digunakan untuk proyek dengan kebutuhan yang stabil dan terdefinisi dengan baik sejak awal, seperti halnya sistem informasi organisasi yang memiliki struktur dan proses bisnis yang jelas \citep{Balaji_2012_Waterfall, Akinsola_2020_SDLC}. Analisis metodologi SDLC kontemporer menegaskan bahwa pemilihan model pengembangan harus disesuaikan dengan karakteristik proyek, di mana \textit{Waterfall} memberikan struktur yang jelas untuk proyek dengan kebutuhan yang telah terdokumentasi dengan baik \citep{Vallon_2018_AgilePracticesGSD, Jabangwe_2018_MobileProcessModelsSLR}.

Menurut \citet{Sommerville_2016_SE}, model \textit{Waterfall} terdiri dari lima tahapan utama yang dilaksanakan secara berurutan:

\begin{enumerate}
  \item \textbf{\textit{Requirement Definition}}: Tahap ini melibatkan pengumpulan dan dokumentasi kebutuhan sistem secara komprehensif melalui wawancara dengan pemangku kepentingan, observasi proses bisnis yang ada, dan analisis dokumen organisasi. Keluaran dari tahap ini adalah dokumen spesifikasi kebutuhan fungsional dan non-fungsional yang menjadi landasan untuk tahap selanjutnya. Dalam konteks sistem informasi INFINITE, tahap ini mencakup identifikasi kebutuhan pengelolaan data keanggotaan, kepengurusan, perlombaan, dan testimoni \citep{Nuseibeh_2000_Requirements}.

  \item \textbf{\textit{System and Software Design}}: Berdasarkan spesifikasi kebutuhan, tahap ini merancang arsitektur sistem secara menyeluruh, termasuk desain basis data, desain antarmuka pengguna, desain API, dan pemilihan teknologi. Perancangan dilakukan pada dua tingkat: \textit{high-level design} yang mendefinisikan arsitektur sistem secara keseluruhan, dan \textit{detailed design} yang merinci setiap komponen dan modul. Untuk sistem INFINITE, tahap ini menghasilkan diagram \textit{Entity-Relationship Diagram} (ERD), \textit{wireframe} antarmuka, spesifikasi \textit{endpoint} API, dan dokumentasi arsitektur \textit{headless} \citep{Bass_2017_SoftwareArchitecture}.

  \item \textbf{\textit{Implementation and Unit Testing}}: Tahap ini melibatkan penerjemahan desain sistem menjadi kode program yang dapat dieksekusi. Pengembang menulis konfigurasi basis data PostgreSQL, kode sumber untuk komponen \textit{backend} menggunakan Laravel, serta komponen \textit{frontend} menggunakan Next.js . Implementasi dilakukan dengan mengikuti praktik terbaik seperti konvensi penamaan yang konsisten, dokumentasi kode, dan penggunaan sistem kontrol versi Git \citep{Foucault_2019_CodingPractices}.

  \item \textbf{\textit{Integration and System Testing}}: Setelah implementasi, sistem diuji secara menyeluruh untuk memastikan bahwa ia berfungsi sesuai dengan spesifikasi kebutuhan dan bebas dari kesalahan. Pengujian dilakukan pada berbagai tingkat: \textit{unit testing} untuk menguji komponen individual, \textit{integration testing} untuk menguji interaksi antar komponen, \textit{system testing} untuk menguji sistem secara keseluruhan, dan \textit{user acceptance testing} untuk validasi oleh pengguna akhir. Dalam penelitian ini, pengujian juga mencakup evaluasi kualitas perangkat lunak berdasarkan standar ISO/IEC 25010:2023 untuk \textit{Product Quality} dan ISO/IEC 25019:2023 untuk \textit{Quality in Use} \citep{Alyahya_2020_CrowdsourcedTestingSLR}.

  \item \textbf{\textit{Operation and Maintenance}}: Tahap akhir melibatkan pengoperasian sistem di lingkungan produksi, pelatihan pengguna, serta pemeliharaan berkelanjutan untuk memperbaiki bug, melakukan pembaruan, dan menambahkan fitur baru sesuai kebutuhan. Otomatisasi \textit{deployment} melalui GitHub Actions dan ArgoCD dengan GitOps memastikan proses rilis yang konsisten dan dapat diandalkan \citep{Ebert_2016_DevOps}.
\end{enumerate}

Keunggulan model \textit{Waterfall} meliputi: (1)~struktur yang jelas dan mudah dipahami, memudahkan manajemen proyek dan estimasi waktu serta biaya; (2)~dokumentasi yang komprehensif pada setiap tahap, memfasilitasi transfer pengetahuan dan pemeliharaan jangka panjang; dan (3)~cocok untuk proyek dengan kebutuhan yang stabil dan tim yang kurang berpengalaman dengan metodologi \textit{Agile}. Namun demikian, model ini juga memiliki keterbatasan, yaitu: (1)~kurang fleksibel terhadap perubahan kebutuhan di tengah proyek; (2)~pengguna baru bisa melihat hasil akhir setelah semua tahap selesai; dan (3)~risiko menemukan masalah desain atau kebutuhan di tahap akhir yang mahal untuk diperbaiki. Meskipun demikian, untuk konteks sistem informasi organisasi mahasiswa dengan kebutuhan yang relatif terdefinisi dengan baik dan sumber daya yang terbatas, model \textit{Waterfall} menawarkan pendekatan yang terstruktur dan dapat dikelola dengan baik \citep{Petersen_2015_WaterfallAgile, Akinsola_2020_SDLC}.

\subsubsection{Pengujian Kualitas Perangkat Lunak}

Pengujian kualitas perangkat lunak adalah proses sistematis untuk mengevaluasi atribut atau kemampuan suatu produk perangkat lunak guna menentukan apakah produk tersebut memenuhi kebutuhan yang ditetapkan dan sesuai dengan standar kualitas yang diharapkan \citep{Alyahya_2020_CrowdsourcedTestingSLR}. Menurut \citet{Pressman_2014_SE}, pengujian kualitas tidak hanya berfokus pada identifikasi kesalahan atau bug, tetapi juga pada verifikasi bahwa perangkat lunak memenuhi spesifikasi kebutuhan fungsional dan non-fungsional, serta validasi bahwa produk akhir sesuai dengan ekspektasi dan kebutuhan pengguna. Pengujian kualitas perangkat lunak memainkan peran kritis dalam memastikan keandalan, keamanan, dan kepuasan pengguna terhadap sistem yang dikembangkan.

International Organization for Standardization (ISO) adalah badan standardisasi internasional independen yang terdiri dari badan-badan standardisasi nasional dari berbagai negara. ISO mengembangkan dan menerbitkan standar internasional yang mencakup berbagai bidang, termasuk teknologi informasi dan rekayasa perangkat lunak. Standar ISO memberikan spesifikasi teknis, pedoman, dan karakteristik terbaik yang memastikan bahwa produk, layanan, dan sistem aman, dapat diandalkan, dan berkualitas tinggi \citep{ISO_2026_About}. Dalam konteks kualitas perangkat lunak, ISO bekerja sama dengan International Electrotechnical Commission (IEC) untuk menghasilkan standar ISO/IEC yang diakui secara global.

ISO/IEC 25010:2023 adalah standar internasional yang merupakan bagian dari rangkaian standar \textit{Systems and Software Quality Requirements and Evaluation} (SQuaRE) \citep{ISO25010_2023}. Standar ini memperbarui model kualitas produk dengan menambahkan karakteristik \textit{safety}, mengganti \textit{usability} dengan \textit{interaction capability}, serta memisahkan model \textit{quality-in-use} ke dalam standar tersendiri ISO/IEC 25019:2023 \citep{ISO25019_2023}. Standar ini menyediakan kerangka kerja yang lebih komprehensif untuk mengevaluasi kualitas produk teknologi informasi dan komunikasi serta perangkat lunak. Penelitian menunjukkan bahwa penerapan ISO/IEC 25010 memerlukan metodologi pengujian yang sistematis untuk mengukur kecukupan informasi dalam penilaian kualitas perangkat lunak \citep{Hovorushchenko_2018_ISO25010}. Model kualitas dalam ISO/IEC 25010:2023 dan ISO/IEC 25019:2023 menjelaskan dua standar model:

\begin{enumerate}
  \item \textbf{\textit{Product Quality}} (Kualitas Produk): Model ini mendefinisikan sembilan karakteristik kualitas yang dapat diukur secara internal maupun eksternal dari produk perangkat lunak. Karakteristik-karakteristik tersebut adalah: (1)~\textit{functional suitability} yang mengukur kelengkapan, keakuratan, dan kesesuaian fungsi; (2)~\textit{performance efficiency} yang mengukur performa dalam hal waktu, sumber daya, dan kapasitas; (3)~\textit{compatibility} yang mengukur kemampuan koeksistensi dan interoperabilitas dengan sistem lain; (4)~\textit{interaction capability} yang mengukur kemampuan interaksi pengguna dengan sistem untuk menyelesaikan tugas; (5)~\textit{reliability} yang mengukur ketersediaan, toleransi kesalahan, dan kemampuan pemulihan; (6)~\textit{security} yang mengukur kerahasiaan, integritas, dan akuntabilitas; (7)~\textit{maintainability} yang mengukur kemudahan modifikasi dan pengujian; (8)~\textit{flexibility} yang mengukur kemampuan adaptasi dan skalabilitas; serta (9)~\textit{safety} yang mengukur kemampuan produk untuk menghindari bahaya terhadap nyawa, kesehatan, dan lingkungan . Standar ISO/IEC 25010 telah diterapkan secara luas sebagai panduan untuk seleksi arsitektur perangkat lunak berdasarkan karakteristik kualitas yang diinginkan \citep{Haoues_2017_ISO25010Architecture}. Penggunaan metrik kualitas yang tepat sangat penting dalam proses evaluasi, dan penelitian terkini memberikan rekomendasi berbasis riset untuk pemilihan dan implementasi metrik kualitas perangkat lunak \citep{Mladenova_2020_QualityMetrics}.

  \item \textbf{\textit{Quality in Use}} (Kualitas dalam Penggunaan): Model ini mengukur efek dan pengaruh pada pemangku kepentingan ketika produk digunakan dalam konteks penggunaan yang spesifik. Model ini terdiri dari tiga karakteristik utama: (1)~\textit{beneficialness} yang mencakup \textit{usability} (efektivitas, efisiensi, dan kepuasan), \textit{accessibility}, dan \textit{suitability}; (2)~\textit{freedom from risk} yang mengukur mitigasi risiko ekonomi, kesehatan, lingkungan, dan risiko terhadap nyawa manusia; serta (3)~\textit{acceptability} yang mengukur keterimaan produk oleh pengguna dan pemangku kepentingan \citep{Wagner_2018_QualityInUse}.
\end{enumerate}

Dalam konteks pengujian sistem informasi organisasi INFINITE, empat karakteristik dipilih sebagai kriteria utama berdasarkan relevansinya dengan tujuan dan sifat produk yang dikembangkan:

\begin{enumerate}
  \item \textbf{\textit{Functional Suitability}} (Kesesuaian Fungsional): Karakteristik ini mengukur sejauh mana perangkat lunak menyediakan fungsi-fungsi yang memenuhi kebutuhan yang dinyatakan dan tersirat ketika digunakan dalam kondisi tertentu. Karakteristik ini mencakup tiga sub-karakteristik: \textit{functional completeness} yang mengukur kelengkapan fungsi untuk mencapai tujuan pengguna, \textit{functional correctness} yang mengukur keakuratan hasil yang disediakan, dan \textit{functional appropriateness} yang mengukur kesesuaian fungsi untuk tugas dan tujuan tertentu \citep{Abana_Sy_2022_ISO25010_Evaluation}.

  \item \textbf{\textit{Performance Efficiency}} (Efisiensi Kinerja): Karakteristik ini mengukur performa sistem dalam hal penggunaan waktu dan sumber daya ketika melakukan fungsi tertentu dalam kondisi yang ditetapkan. Sub-karakteristik meliputi \textit{time behavior} yang mengukur waktu respons dan pemrosesan, \textit{resource utilization} yang mengukur penggunaan CPU, memori, dan bandwidth, serta \textit{capacity} yang mengukur batas maksimum parameter sistem \citep{Kurniasari_Rochimah_2025_PerformanceEfficiency}.

  \item \textbf{\textit{Interaction Capability}} (Kemampuan Interaksi): Karakteristik ini mengukur sejauh mana produk atau sistem dapat diinteraksi oleh pengguna tertentu untuk bertukar informasi antara pengguna dan produk atau sistem, melalui antarmuka pengguna, termasuk masukan dari pengguna dan keluaran dari produk atau sistem, untuk tujuan menyelesaikan tugas dalam berbagai konteks penggunaan. Sub-karakteristik meliputi \textit{appropriateness recognizability}, \textit{learnability}, \textit{operability}, \textit{user error protection}, \textit{user engagement}, \textit{inclusivity}, \textit{user assistance}, dan \textit{self-descriptiveness} \citep{Cayola_Macias_2018_UsabilityMethods}.

  \item \textbf{\textit{Beneficialness}} (Kemanfaatan): Karakteristik ini merupakan bagian dari model \textit{quality-in-use} dalam ISO/IEC 25019:2023, yang mengukur sejauh mana manfaat yang dihasilkan dari penggunaan produk atau sistem bagi pemangku kepentingan dalam konteks penggunaan yang spesifik. Karakteristik ini mencakup sub-karakteristik \textit{usability} yang meliputi efektivitas, efisiensi, dan kepuasan; \textit{accessibility}; serta \textit{suitability}.
\end{enumerate}

\section{Hasil Penelitian yang Relevan}

Berdasarkan kajian literatur yang telah dilakukan, terdapat beberapa hasil penelitian yang relevan dengan topik penelitian ini, di antaranya:

\begin{enumerate}
  \item Penelitian oleh \citet{Pramitasari_Nurgiyatna_2019_UKMMarchingBand} berjudul \textit{Sistem Informasi Unit Kegiatan Mahasiswa Marching Band Universitas Muhammadiyah Surakarta berbasis Web} bertujuan mengembangkan sistem informasi UKM berbasis web untuk mendukung penyebaran informasi, pengelolaan data kegiatan, serta inventaris organisasi. Penelitian tersebut menggunakan metode \textit{Waterfall} dengan implementasi PHP dan MySQL. Hasil pengujian \textit{black box} menunjukkan fitur sistem berjalan sesuai kebutuhan, dan evaluasi pengguna (37 responden) menunjukkan tingkat persetujuan sebesar 91,1\% bahwa sistem membantu pengelolaan data dan penyebaran informasi UKM. Relevansi terhadap penelitian ini terletak pada kesamaan konteks organisasi kemahasiswaan dan tujuan digitalisasi proses administrasi organisasi.

  \item Penelitian oleh \citet{Widjaja_Prasojo_2022_UKMKarangturi} berjudul \textit{Perancangan Sistem Informasi Unit Kegiatan Mahasiswa Universitas Nasional Karangturi Berbasis Web} berfokus pada perancangan sistem informasi UKM untuk mengoptimalkan pengelolaan aktivitas kemahasiswaan. Metode yang digunakan meliputi wawancara dan observasi untuk pengumpulan data, analisis PIECES, perancangan UML, serta pengembangan perangkat lunak dengan model \textit{Waterfall}. Hasil penelitian menunjukkan sistem yang dirancang memberikan manfaat bagi mahasiswa, pengurus UKM, dan admin dalam menjalankan aktivitas organisasi. Penelitian ini relevan karena menegaskan bahwa pendekatan \textit{Waterfall} dan pemodelan terstruktur efektif untuk kebutuhan sistem informasi UKM di lingkungan perguruan tinggi.

  \item Penelitian oleh \citet{Tofani_Sukya_2023_UKMEnglishClub} berjudul \textit{Sistem Informasi Manajemen Kegiatan UKM English Club PSDKU Polinema di Kediri Berbasis Framework Laravel} mengembangkan sistem manajemen kegiatan UKM yang sebelumnya dikelola secara manual dan tidak terintegrasi. Sistem dibangun menggunakan PHP, Laravel, dan PostgreSQL dengan empat peran pengguna (DPK, ketua UKM, sekretaris UKM, dan admin kegiatan). Fitur utama mencakup pengelolaan program kerja, persetujuan program kerja, pengelolaan proposal kegiatan, serta LPJ kegiatan. Hasil pengujian menunjukkan seluruh fungsi utama berjalan sesuai harapan. Relevansi penelitian ini terhadap penelitian yang dilakukan terletak pada kesamaan penggunaan teknologi Laravel dan PostgreSQL serta fokus pada integrasi proses manajemen kegiatan organisasi mahasiswa. Namun, penelitian tersebut belum membahas pembaruan implementasi menuju \textit{deployment} berbasis kontainer.

  \item Penelitian oleh \citet{Feriawan_Paramitha_Dewi_2024_TAKPrimakara} berjudul \textit{Rancang Bangun Sistem Informasi Manajemen Transkrip Aktivitas Kemahasiswaan Primakara University Menggunakan Metode Extreme Programming} membahas pengembangan sistem informasi aktivitas kemahasiswaan berbasis web untuk menggantikan proses manual berbasis email yang tidak efisien. Sistem dikembangkan menggunakan Next.js dan JavaScript dengan metode \textit{Extreme Programming} (XP), melalui tahapan perencanaan, perancangan UML, implementasi, dan pengujian \textit{black box}. Hasil penelitian menunjukkan sistem memenuhi kebutuhan fungsional dan mendukung pemantauan poin aktivitas secara \textit{real-time}. Relevansi penelitian ini terletak pada kesamaan domain sistem informasi kemahasiswaan dan penggunaan teknologi \textit{frontend} modern, sekaligus memberikan pembanding metodologis terhadap model \textit{Waterfall} yang digunakan pada penelitian ini. Meskipun demikian, penelitian tersebut belum menekankan orkestrasi layanan dan otomasi rilis berbasis kontainer.
\end{enumerate}

Berdasarkan keempat penelitian tersebut, dapat disimpulkan bahwa pengembangan sistem informasi organisasi kemahasiswaan berbasis web secara konsisten mampu meningkatkan efektivitas pengelolaan data, transparansi proses administrasi, dan kemudahan akses informasi bagi pemangku kepentingan. Namun, penelitian terdahulu umumnya masih berfokus pada pengembangan fitur dan pengujian fungsional, serta belum menempatkan strategi \textit{deployment} modern sebagai fokus utama. Posisi penelitian ini berada pada penguatan integrasi pengelolaan organisasi INFINITE dengan arsitektur \textit{headless} (Laravel sebagai \textit{backend} dan Next.js sebagai \textit{frontend}), pembaruan implementasi melalui kontainer dan Kubernetes berbasis GitOps, serta evaluasi kualitas perangkat lunak berbasis ISO/IEC 25010:2023 untuk \textit{Product Quality} dan ISO/IEC 25019:2023 untuk \textit{Quality in Use} pada karakteristik yang relevan dengan kebutuhan pengguna.

\section{Kerangka Pikir}

Berikut ini adalah kerangka pikir yang mendasari penelitian pengembangan sistem informasi organisasi INFINITE dengan pendekatan Laravel \textit{headless}, Next.js, dan strategi \textit{deployment} berbasis kontainer dan Kubernetes:

\begin{figure}[H]
  \centering
  \caption{Diagram Kerangka Pikir Penelitian}
  \label{fig:kerangka-pikir}
  \begin{tikzpicture}[
      node distance=0.8cm,
      widebox/.style={
          rectangle,
          draw=black,
          minimum width=13cm,
          text width=12.4cm,
          align=center,
          inner sep=10pt,
          font=\small\linespread{1}\selectfont
        },
      arrow/.style={->, >=stealth, thick}
    ]

    \node (problem) [widebox] {\textbf{Masalah}\\
      \begin{minipage}{12.4cm}
        \medskip
        \begin{enumerate}
          \item Sistem informasi yang ada terpecah dalam dua sistem berupa sistem Laravel di \textit{shared hosting} dan sistem Strapi di server terpisah, menciptakan fragmentasi dan potensi inkonsistensi data.
          \item Sistem informasi Strapi yang ada tidak mendukung integrasi dengan protokol OIDC untuk memungkinkan autentikasi terpusat dengan sistem SSO.
          \item Sistem informasi Laravel yang ada tidak memiliki \textit{image} kontainer untuk dapat di-\textit{deploy} pada platform Kubernetes.
          \item Pengoperasian dua infrastruktur secara bersamaan (\textit{shared hosting} dan VPS) mengakibatkan inefisiensi biaya operasional.
          \item Sistem informasi Laravel yang ada menggunakan arsitektur \textit{monolithic} yang tidak menyediakan API untuk integrasi dengan aplikasi lain.
        \end{enumerate}
      \end{minipage}};

    \node (solution) [widebox, below=of problem] {\textbf{Solusi}\\Pengembangan sistem informasi organisasi INFINITE dengan pendekatan \textit{headless} menggunakan Laravel sebagai \textit{backend} dan Next.js sebagai \textit{frontend}, integrasi autentikasi SSO dengan OIDC, serta implementasi strategi \textit{deployment} berbasis kontainer dan Kubernetes dengan pendekatan GitOps.};

    \node (dev) [widebox, below=of solution] {\textbf{Implementasi Pengembangan Perangkat Lunak (\textit{Waterfall})}\\Melalui tahap \textit{Requirements Definition}, \textit{System and Software Design}, \textit{Implementation and Unit Testing}, \textit{Integration and System Testing}, serta \textit{Operation and Maintenance}.};

    \node (test) [widebox, below=of dev] {\textbf{Pengujian (ISO/IEC 25010:2023 dan ISO/IEC 25019:2023)}\\\textit{Functional Suitability}, \textit{Performance Efficiency}, \textit{Interaction Capability}, dan \textit{Beneficialness}.};

    \node (result) [widebox, below=of test] {\textbf{Hasil}\\Sistem informasi organisasi INFINITE dengan Laravel \textit{headless} dan Next.js yang terintegrasi dengan SSO OIDC, serta di-\textit{deploy} menggunakan kontainer dan Kubernetes berbasis GitOps.};

    \draw [arrow] (problem) -- (solution);
    \draw [arrow] (solution) -- (dev);
    \draw [arrow] (dev) -- (test);
    \draw [arrow] (test) -- (result);
  \end{tikzpicture}
\end{figure}

\section{Pertanyaan Penelitian}

\begin{enumerate}
  \item Bagaimana tahapan \textit{Requirements Definition} untuk mengembangkan sistem informasi organisasi INFINITE dengan pendekatan \textit{headless} menggunakan Laravel dan Next.js yang terintegrasi autentikasi SSO dengan OIDC?
  \item Bagaimana tahapan \textit{System and Software Design} untuk mengembangkan sistem informasi organisasi INFINITE dengan pendekatan \textit{headless} menggunakan Laravel dan Next.js yang terintegrasi autentikasi SSO dengan OIDC?
  \item Bagaimana tahapan \textit{Implementation and Unit Testing} untuk mengembangkan sistem informasi organisasi INFINITE dengan pendekatan \textit{headless} menggunakan Laravel dan Next.js yang terintegrasi autentikasi SSO dengan OIDC?
  \item Bagaimana tahapan \textit{Integration and System Testing} untuk mengembangkan sistem informasi organisasi INFINITE dengan pendekatan \textit{headless} menggunakan Laravel dan Next.js yang terintegrasi autentikasi SSO dengan OIDC?
  \item Bagaimana hasil pengujian kualitas sistem informasi organisasi INFINITE berdasarkan aspek \textit{functional suitability}?
  \item Bagaimana hasil pengujian kualitas sistem informasi organisasi INFINITE berdasarkan aspek \textit{performance efficiency}?
  \item Bagaimana hasil pengujian kualitas sistem informasi organisasi INFINITE berdasarkan aspek \textit{usability}?
  \item Bagaimana hasil pengujian kualitas sistem informasi organisasi INFINITE berdasarkan aspek \textit{beneficialness}?
\end{enumerate}
